\section{Cálculo del coeficiente}\label{sec:numeval}
La fórmula~\eqref{eq:main} permite calcular $\cerr$ a partir del perfil físico. Para cada valor de $x$
hay que sumar las dieciséis contribuciones a $d(x)$ sobre las posiciones del tablero. Las primeras
evaluaciones, que promediaban unas pocas casillas centrales, daban errores espaciales de $10^{-5}$
a $10^{-3}$. En particular, había que volver a evaluar la resolvente de líneas para configuraciones
que compartían una casilla. El método de rectángulos suma todos los
centros de un patrón de incidencia en una misma malla. Una vez obtenidos los valores en $x$, se
integra en esa variable y se añaden los términos de frontera.

\subsection{El método de rectángulos}\label{sec:rectangulos}
La idea procede de la técnica de rectángulos del término constante~\cite[\S4]{tc}. Allí la norma de
una matriz de error $E$ se controla con
\[
 \operatorname{tr}\bigl((EE^{\mathsf T})^2\bigr)=\sum_{i,i',j,j'}E_{ij}E_{i'j}E_{i'j'}E_{ij'},
\]
una suma sobre rectángulos de dos filas y dos columnas. Cada sumando se desarrolla en dieciséis
patrones $T\subseteq\{1,2,3,4\}$, según qué vértices llevan sus líneas diagonales, con signo
$(-1)^{4-|T|}$. Como $1-4+6-4+1=0$, el orden principal se cancela exactamente. El método lleva esa
contabilidad a las dieciséis contribuciones de $d(x)$. Cada una es una suma sobre caminos cerrados
cortos de casillas unidas por líneas, triángulos y caminos de hasta cuatro vértices con las cuatro
familias, y los patrones se agrupan antes de integrar. Así cada grupo es la traza de un producto de
matrices pequeñas en una misma malla.

La precisión necesaria viene del extremo $x=0$, donde $\numerr/x^2$ amplifica los errores. Con una regla de
Gauss--Legendre de dieciséis puntos en $[0,1]$, el error de una fila en el primer nodo entra en la
integral multiplicado por $w_1/x_1^2\approx483$. Por eso cerca del origen se usa una serie. En el
cálculo numérico se busca una precisión del orden de $10^{-10}$ en los nodos y se evalúan también
las derivadas en el origen. La reordenación de las sumas es exacta en cada malla. Por separado,
con un perfil de malla de marginales exactas se midieron los residuos de las dos cancelaciones en
punto flotante. Para $5\le m\le48$, en ambas paridades, los valores observados cumplieron
$|\numerr_m(0)|\le7{,}5\cdot10^{-13}$ y $|d_m'(0)|\le1{,}3\cdot10^{-11}$ \status{N}.

\subsection{Una malla cerrada bajo las intersecciones de líneas}
Numeramos las filas y columnas desde cero y situamos la casilla $(i,j)$ en
$((i+1/2)/m,(j+1/2)/m)$. Hay $6m-2$ líneas. Las diagonales con índices enteros $i-j$ e $i+j$
que pasan por la casilla tienen la misma paridad. Recíprocamente, dos índices diagonales de igual
paridad dan una casilla si su intersección cae dentro del tablero. En el continuo las dos
paridades dan cuatro hojas oblicuas, además de las hojas axiales $R,C$. Sea $\nu_1$ la densidad
del perfil de una solución completa.
Para $x>0$ tomamos $U_f=u_f+x(1-u_f)$ sobre una línea oblicua y
$W=\nu_1/(U_DU_A)$. Cuando las líneas $a,b$ se cortan, $K_{ab}$ es el valor de $\nu_1$ en la
intersección multiplicado por $x/U_f$ por cada línea oblicua $f$ de la pareja. Las demás entradas se
anulan. La matriz diagonal $\rho$ vale $1/m$ en las axiales y $U_f/(xm)$ en las oblicuas. Sea $v$ el
vector que vale $1,-1,0$ en filas, columnas y oblicuas, respectivamente. El núcleo de las contracciones
es
\[
 \Pi=vv^{\mathsf T}/2,\qquad
 L_0=K(I+\rho K+\rho\Pi)^{-1}\rho K-\tfrac32vv^{\mathsf T}.
\]
El integrando en $x=0$ se define por continuación. En la malla se conservan ambas hojas cuando un
observable toma valores en una sola paridad.

Para un peso de casilla $h$, definimos $N_h^\circ(a,b)=h(a\cap b)$ si las líneas distintas
$a,b$ se cortan en una casilla y cero en las demás entradas. Con la incidencia casilla--línea
$\mathcal I_{ea}=\mathbf 1_{\{e\in a\}}$, la matriz completa es
\[
 N_h^{\rm inc}=\mathcal I^{\mathsf T}\operatorname{diag}(h)\mathcal I
              =N_h^\circ+\operatorname{diag}(\mathcal I^{\mathsf T}h).
\]
Su entrada $(a,a)$ suma $h$ sobre la línea $a$. Los productos matriciales suman recorridos por
incidencias de líneas. La traza identifica el último índice con el primero y las matrices diagonales
$D_i$ insertan las medidas de línea. Por ejemplo, $\operatorname{tr}(N_{h_1}^\circ D_1
N_{h_2}^\circ D_2 N_{h_3}^\circ D_3)$ suma los productos de los tres pesos sobre los triángulos.
La Figura~\ref{fig:rectangulos} muestra dos rectángulos de la parte de
$\operatorname{tr}(N_{h_1}^\circ N_{h_2}^\circ N_{h_3}^\circ N_{h_4}^\circ)$ que alterna filas y columnas.
Cada vértice de esos rectángulos cruza dos líneas distintas. La traza general de $f_6$ del
Lema~\ref{lem:f6} usa $N_h^{\rm inc}$ e incluye recorridos de hasta cuatro vértices.
Los dieciséis aportes a $d$ se evalúan mediante productos matriciales, formas bilineales y sumas
locales. La traza reordena una suma finita de manera exacta. Compararla con un recuento directo en
la misma malla comprueba la implementación. Para estimar la integral del continuo hay que variar $m$.
\begin{figure}[htbp]
\centering
\begin{tikzpicture}[x=.56cm,y=.56cm,
  vertice/.style={circle,draw=btinta,fill=white,text=btinta,font=\scriptsize\bfseries,
    inner sep=0pt,minimum size=3.5mm},
  paso/.style={->,>=stealth,draw=btinta,line width=.85pt}]
  \foreach \dx/\izq/\der in {0/1/3,9.7/2/4}{
    \begin{scope}[shift={(\dx,0)}]
      \bcuadriculacero{6}
      \node[vertice] (v1) at ({\izq+.5},4.5) {1};
      \node[vertice] (v2) at ({\izq+.5},1.5) {2};
      \node[vertice] (v3) at ({\der+.5},1.5) {3};
      \node[vertice] (v4) at ({\der+.5},4.5) {4};
      \draw[paso] (v1)--(v2);
      \draw[paso] (v2)--(v3);
      \draw[paso] (v3)--(v4);
      \draw[paso] (v4)--(v1);
      \node[font=\small] at (3,-.73)
        {$R_1\to C_{\izq}\to R_4\to C_{\der}\to R_1$};
    \end{scope}
  }
\end{tikzpicture}
\caption{Dos sumandos de una traza de cuatro incidencias en la malla $6\times6$.
El vértice $k$ marca la intersección donde $N_{h_k}^{\circ}$ evalúa $h_k$. A la izquierda, el producto es
$h_1(1,1)h_2(4,1)h_3(4,3)h_4(1,3)$. A la derecha, cada columna aumenta en uno.
El cuarto factor vuelve a la fila inicial $R_1$. La suma recorre índices de línea ordenados.
Cada fórmula fija las medidas y las multiplicidades.}\label{fig:rectangulos}
\end{figure}


\subsection{De la suma de malla al valor en cada nodo}
En las coincidencias accidentales se conserva el valor límite del continuo. En la suma de cuadrados
se asigna peso medio al lado de longitud cero. Al cambiar $m$ cambia también la posición de los
pliegues respecto de los nodos. Por eso se toman mallas de una misma clase de congruencia para
cada ajuste en potencias pares de $1/m$, como en la extrapolación de
Richardson~\cite{richardson}. Si hay suficientes mallas nuevas para $Q_2$, se ajustan por
separado las clases $0$ y $12$ módulo
$24$, con potencias $2,4,6$. Se toma la media de las dos ordenadas en el origen. Su
diferencia, un cambio de la malla mínima y el mayor residuo del ajuste miden la sensibilidad. Las
barras resultantes tienen estado \status{N}. Certificar esta extrapolación exigiría acotar el resto espacial
en cada clase.

Para la suma de frontera el análisis numérico da
\[
 J=-8{,}2386091222\ \pm5{,}0\cdot10^{-9},\qquad
 P_0=-7{,}7933777386\ \pm5{,}0\cdot10^{-9}\qquad\text{\status{N}}.
\]
Para $J$ se suma una orientación del par y se incluye la otra por simetría. Al añadir
$P_T=0{,}4452313836$ se obtiene $P_0$. La interacción del cuadrado entra por separado con signo
negativo, con $S=15{,}6334582804$. La Sección~\ref{sec:numcert} da los encierros rigurosos
de $P_T$ y $S$.

\subsection{El extremo singular y el presupuesto de error}\label{sec:presupuesto}
Para integrar con el peso $x^{-2}$ conviene extraer primero la parte que tiene una primitiva
explícita. Separamos el funcional de torres mediante
\[
 \numerr_F=F-F(0)-xF',\qquad
 \numerr_6=\numerr-\numerr_F=d+\mathcal V-x\mathcal V'+F(0)-1/12.
\]
La identidad $\numerr_F/x^2=-[(F(x)-F(0))/x]'$ da la integral exacta $F(0)+F'(0)-F(1)$ en $[0,1]$.
Para $\numerr_6/x^2$, los términos constante y lineal de $\numerr_6$ se anulan. El
Apéndice~\ref{app:numquad} acota el error al integrar esta parte con una serie obtenida en
$|x|=0{,}03$ y dos reglas de Gauss--Legendre, bajo las tolerancias de muestreo allí indicadas.

Para obtener el valor numérico de $\cerr$, la réplica forma la serie del $\numerr$ completo a partir de
$d$ y $\mathcal V+F$ en los 32 puntos del círculo (17 independientes por conjugación).
Evalúa $\numerr/x^2$ en seis nodos de $[0{,}03;0{,}06]$ y veinte de $[0{,}06;1]$. Como cotejo calcula
también la contribución de $\numerr_F/x^2$ con la misma serie y ambas reglas. La fórmula de los extremos,
evaluada con los datos de la réplica, menos esa integral numérica da $-3{,}77\cdot10^{-13}$ \status{N}.
La tabla da el presupuesto empírico \status{N} de la evaluación. Para los nodos y $P_0$ se usa la
sensibilidad al refinamiento de la malla. Las dos últimas filas indican el alcance de sus cotas.

\begin{center}\small
\begin{tabular}{lr}
Origen & Presupuesto empírico \status{N}\\ \hline
Datos nodales del panel de veinte puntos & $2{,}8\cdot10^{-8}$\\
Datos del círculo de radio $0{,}03$ & $1{,}4\cdot10^{-8}$\\
Valor de frontera $P_0$ & $5{,}0\cdot10^{-9}$\\
Datos nodales del panel de seis puntos & $2{,}3\cdot10^{-9}$\\
Serie y cuadratura de $\numerr_6/x^2$ \status{D si C} & $8{,}74\cdot10^{-10}$\\
Error absoluto del valor de $S$ usado \status{O} & $7{,}044\cdot10^{-10}$
\end{tabular}
\end{center}
Para $S=15{,}6334582804$ se ha tomado la mayor distancia a los extremos certificados de la
Sección~\ref{sec:numcert}. El semiancho certificado de $P_T$ es aproximadamente
$9{,}1\cdot10^{-4}$ \status{O}. La fila de $P_0=P_T+J$ mide la variación observada en malla \status{N}.

El cálculo da $0{,}4055003786$. Las mayores contribuciones al error proceden de los valores en el
panel de veinte puntos y en el círculo. La combinación cuadrática de los indicadores es
$3{,}2\cdot10^{-8}$ y la suma lineal es $5{,}1\cdot10^{-8}$ con la cota de $S$.
Al cambiar el ajuste de $Q_2$, el resultado varió
entre $-3{,}7\cdot10^{-9}$ y $+7{,}6\cdot10^{-9}$ respecto de la referencia. Con estos controles empíricos
se comunica $0{,}40550038\pm3\cdot10^{-8}$ \status{N}, con la estimación conservadora
$0{,}4055004\pm2\cdot10^{-7}$ \status{N}.

Redondeados a ocho decimales, los términos de~\eqref{eq:main} en esta evaluación son los siguientes
\status{N}. La integral se da como el total menos los demás términos, y la suma lleva una cancelación
de casi dos órdenes de magnitud.
\begin{center}\small\begin{tabular}{lr}
Término & Valor\\ \hline
$F(1)-\tfrac1{12}$ & $-0{,}08333333$\\
$\int_0^1\numerr(x)x^{-2}\,dx$ & $23{,}91566973$\\
$P_T$ & $0{,}44523138$\\
$J=-B_2+L_2$ & $-8{,}23860912$\\
$-S$ & $-15{,}63345828$\\ \hline
$\cerr$ & $0{,}40550038$
\end{tabular}\end{center}
